\section{Calibration and the exact boundary resolvent}\label{sec:boundary}
\subsection{A common calibration amplitude on both parities}
The earlier exact identities give
\[
 Z_{2m}(c)=2^m\prod_{k=1}^m\frac{A_{2k}}{A_{2k-1}},\qquad
 Z_n(c)=c2^{n-1}A_n/Z_{n-1}(c)\quad(n\text{ odd}).
\]
Let $O_m$ denote the number of off-diagonally symmetric ASMs of
order $2m$. The classical products in \cite{bfk} are
\[
 A_N=\prod_{i=0}^{N-1}\frac{(3i+1)!}{(N+i)!},\qquad
 O_m=\prod_{i=1}^{m}\frac{(6i-2)!}{(2m+2i)!}.
\]
Successive ratios give
\[
 \frac{O_m}{O_{m-1}}
 =\frac{(6m-2)!(2m)!}{(4m)!(4m-2)!}
 =\frac12\frac{A_{2m}}{A_{2m-1}}.
\]
For example, the old denominator factors in the $O_m$ quotient
telescope to $(4m-2)!/(2m)!$. The initial products equal one,
so $Z_{2m}(c)=2^{2m}O_m$ exactly.

The published product asymptotics \cite[Equations (11.3)--(11.6)]{bfk},
using the even-order OSASM branch, state
\begin{align*}
 A_n&\sim C_{\rm ASM}e^{2\alpha n^2}n^{-2\kappa},\\
 2^nO_{n/2}&\sim C_{\rm cal}e^{\alpha n^2+bn}n^{-\kappa}
                    \quad(n\text{ even}),\\
 C_{\rm ASM}&=
 \frac{2^{5/12}\pi^{1/3}e^{\zeta'(-1)/3}}
      {3^{7/36}\Gamma(1/3)^{2/3}}.
\end{align*}
These are known product asymptotics, not an assumption about DSASM
amplitudes. Substitution in the odd exact identity gives the odd
constant
\[
 C_{\rm odd}=(c/2)e^{b-\alpha}C_{\rm ASM}/C_{\rm cal}.
\]
Now $b-\alpha=\tfrac12\log2$ and
$C_{\rm ASM}/C_{\rm cal}^2=\sqrt{2/3}$ by the displayed constants.
Therefore $C_{\rm odd}=C_{\rm cal}$, proving the full-sequence
calibration
\begin{equation}\label{calibrationequiv}
 Z_n(c)\sim C_{\rm cal}e^{\alpha n^2+bn}n^{-\kappa}.
\end{equation}
Only this equivalent is needed below; no finite-size correction is
transferred to a noncalibrated fugacity.

\subsection{The fixed operator, basis and source}
Put $h=\sqrt3$, $a=\pi/(3h)$, $p(y)=y^2+3/4$, and
\begin{equation}\label{realdefinitions}
 w_y(y)=\frac{p(y)}{2(2\cosh(2ay)-1)},\quad
 m(y)=\frac1{1+e^{2ay}},\quad
 \varphi(u)=\frac1{3(1+2\cosh(2au))}.
\end{equation}
Let $\psi_j$ be the real orthonormal polynomials in
$\mathcal H=L^2(w_y\,dy)$ whose leading coefficients are
$(-1)^j/\sqrt{H_j}$. The degree-$<n$ projection is $P_n$.
On polynomials put
\[
 (Tg)(y)=\int_\R e^{au}\varphi(u)g(y+u)\,du,\qquad \mathcal B=mT.
\]
The boundedness of $\mathcal B$ is part of the earlier proof.
The unweighted transform of $\mathcal B$ has kernel
\begin{equation}\label{kernelreal}
 \widetilde{\mathcal B}(y,x)=g_0(y)\varphi(x-y)h_0(x),\quad
 g_0=\sqrt{w_y}\,m e^{-ay},\quad h_0=e^{ay}/\sqrt{w_y},
 \quad C_0(y,x)=\one_{y<0}\varphi(x-y)\one_{x<0}.
\end{equation}
Here $\mathcal E=\widetilde{\mathcal B}-C_0$ is the operator in
\eqref{Arel}. The polynomial inverse of $T$ is
\[
 (\mathcal Kg)(y)=g(y)+e^{i\pi/3}g(y+ih)+e^{-i\pi/3}g(y-ih).
\]
The exact degree-space determinant is
$D_n(t)/D_n(3)=\det(I+(t-3)P_n\mathcal B P_n)$.

Define $d_j=\sqrt{(j+1)(j+2)}$, $h_j=j!(3)_j$, and
\begin{equation}\label{Hnorm}
 H_j=2(27/16)^j
       \frac{j!(3)_j(5/3)_j(7/3)_j}{(3/2)_j(5/2)_j}.
\end{equation}
The affine source $q_t\in\mathcal H$ is
\begin{equation}\label{qsource}
 q_t(y)=\frac{A_t e^{2ay}+(-y-t+3/2)e^{ay}/\sqrt2+C_t
           +[(1-t)y/2-t+3/2]e^{-ay}/\sqrt2}{p(y)\cosh(ay)},
\end{equation}
where $A_t=\sqrt{3/2}(t+1)/4$ and $C_t=\sqrt{3/2}(t-2)$.
Write $q_t=q_3+(t-3)q^{(1)}$ and
\begin{equation}\label{gamma}
 \gamma_j(t)=\left\langle\psi_j,
   (I+(t-3)P_{j+1}\mathcal B^*P_{j+1})^{-1}P_{j+1}q_t\right\rangle,
 \qquad a_j^{\rm src}=\langle\psi_j,q_3\rangle.
\end{equation}
The inner product is linear in its second argument. The inverse in
\eqref{gamma} acts on the finite-dimensional range, never on an
unjustifiably substituted infinite section.

\begin{proposition}[Exact all-parity boundary formula]\label{prop:boundary}
For every $n\ge1$ and $s>0$, with $t=s^2$,
\begin{equation}\label{boundaryexact}
 \frac{sZ_n(s)}{Z_{n+1}(s)}
 =(-1)^{n-1}d_{n-1}\sqrt{h_{n-1}/H_{n-1}}\,\gamma_{n-1}(t).
\end{equation}
It follows that $F_n(s)=\gamma_{n-1}(s^2)/a_{n-1}^{\rm src}$, and
\begin{equation}\label{calcoefficient}
 a_j^{\rm src}\sim\frac{(-1)^j\sqrt2\,3^{-1/4}}{j+1}.
\end{equation}
\end{proposition}
\begin{proof}
We supply the borders and source normalization explicitly.
Let
\[
 K_t(u,v)=\frac{v-u}{1-uv}
       \left(t+\frac{(1+u)(1+v)}{1-u-v}\right),\qquad
 M_{ij}=[u^iv^j]K_t(u,v).
\]
The source Pfaffian formula \cite[Equation (4.1)]{bfk}
and the earlier adjacent determinant give
$H_n=(M_{i,j+1})_{0\le i,j<n}$ with $\det H_n=D_n(t)$.
Set $b_i=M_{i0}$ and $v_t=te_0-b$; directly
$b_0=0$, $b_1=-(t+1)$ and $b_i=-2$ for $i\ge2$.
For even $n$, replacing the last column of $H_n$ by $b$ and cycling
it to the front gives
$-\det(H_n)e_{n-1}^{\mathsf T}H_n^{-1}b=Z_n(s)^2$.
The $e_0$ cofactor vanishes because its complementary skew matrix
has odd size. For odd $n$ the full skew determinant vanishes, while
the $e_0$ complementary skew minor indexed $1,\ldots,n-1$ has
determinant $(Z_n(s)/s)^2$. Including the cofactor signs in either
case proves
\begin{equation}\label{cofactor}
 sZ_n(s)/Z_{n+1}(s)=e_{n-1}^{\mathsf T}H_n^{-1}v_t.
\end{equation}

Use the finite matrices $S$ (lower shift), $L_{ij}=\binom ij$,
$E_{ii}=(-1)^i$, $\Pi_{ij}=\binom{i+j}{i}$, and $T_n(g)=g(S)$.
The earlier finite Pascal factorization, or direct multiplication
of the displayed generating functions, is
\[
 H_n=T_n(B_0)LEJ_n(t)EL^{\mathsf T},\quad
 J_n(t)=I+T_n(G_t)\Pi,
\]
where
\[
 B_0(z)=\frac{(z-2)(z+1)(2z-1)}{(1-z)(z^2-z+1)},\quad
 G_t(z)=\frac{(t-4)z^2+(4-t)z+t-1}{(2-z)(1-z)^2(1+z)}.
\]
The last row of $L^{-\mathsf T}$ is $e_{n-1}^{\mathsf T}$.
The generating function of $v_t$ is $V_t(z)=t+(t+1)z+2z^2/(1-z)$,
so the binomial transform
$(1-z)^{-1}[V_t/B_0](-z/(1-z))$ simplifies to
\[
 w_t(z)=\frac{(t-z)(z^2-z+1)}{(2-z)(1-z)^2(1+z)}.
\]
Consequently \eqref{cofactor} equals
$(-1)^{n-1}e_{n-1}^{\mathsf T}J_n(t)^{-1}w_t^{(n)}$,
where the vector contains its first $n$ coefficients.

Put $D=\diag(d_j)$, $A=D^{-1}SD$ and
$C=D^{-1}VD$ with $V_{ij}=\binom{i+j+2}{i}$.
For $r(z)=1-z+z^2$, $q(z)=2+z-z^2$, the exact finite identity is
$D^{-1}J_nD=I+C+(t-3)[r(A)/q(A)]C$.
Multiplying by $q(A)/r(A)$ gives the old row-action Gram matrix
$F_n^\phi=\mathcal K_n^{\rm row}(I+C)+(t-3)C$.
The forcing is especially simple:
\[
 (q/r)w_t=(t-z)/(1-z)^2,\qquad
 p_i(t)=[t(i+1)-i]/d_i.
\]
Thus \eqref{cofactor} becomes
$(-1)^{n-1}d_{n-1}e_{n-1}^{\mathsf T}(F_n^\phi)^{-1}p(t)$.

Let $\phi_j=u_j/\sqrt{h_j}$ be the Meixner--Pollaczek basis from
\cite[Section 6]{r124}, with
\[
 \sum_{j\ge0}u_j(y)z^j/j!=r(z)^{-3/2}e^{y\eta(z)},\qquad
 \eta'(z)=-1/r(z),\quad\eta(0)=0.
\]
The lower-triangular basis change $\psi=U_n\phi$ has last diagonal
entry $\sqrt{h_{n-1}/H_{n-1}}$. Congruence gives
$F_n^\psi=U_nF_n^\phi U_n^{\mathsf T}=\mathcal K_n^{\rm row}
+(t-3)M_n$, where $(M_n)_{ij}=\langle\psi_i,m\psi_j\rangle$.
The matrix $(\mathcal K_n^{\rm row})^{-1}M_n$ is exactly the matrix
of $P_n\mathcal B^*P_n$ in column coordinates. If
$\ell_t(\phi_i)=p_i(t)$, its transformed source has entries
$\ell_t(T\psi_i)$. We now identify their genuine representing vector.

Define the exponentially decaying density
\[
 \nu_t(y)=\frac{(y+t-3/2)e^{ay}+[(1-t)y+(3-t)/2]e^{-ay}}
                    {2\sqrt2\cosh(3ay)}.
\]
The integral
$\int e^{\eta y}/[2\cosh(3ay)]\,dy=(h/2)\sec(h\eta/2)$
and its derivative yield, for $x=h\eta/2$,
\begin{equation}\label{Lsource}
 L_t(\eta):=\int e^{\eta y}\nu_t(y)\,dy
 =\frac3{\sqrt2}\frac{3t\cos x+\sqrt3(2-t)\sin x}
                            {[1+2\cos(2x)]^2},\quad |\Re\eta|<2a.
\end{equation}
The inverse change is
$z=-2\sin x/(\sqrt3\cos x-\sin x)$ and
$r(z)=3/(\sqrt3\cos x-\sin x)^2$.
Substitution gives
$r(z)^{-3/2}L_t(\eta(z))=(t-z)/[\sqrt2(1-z)^2]$.
Since $\sqrt{h_i}=i!d_i/\sqrt2$, coefficient comparison proves
$\ell_t(g)=\int\nu_tg$ on polynomials.

Set
$\rho_t(y)=\int e^{au}\varphi(u)\nu_t(y-u)\,du$.
Exponential tails justify Fubini for each polynomial and show
$\ell_t(T\psi_i)=\int\psi_i\rho_t$.
Its Laplace transform on $-2a<\Re\eta<a$ is
\[
 \int e^{\eta y}\rho_t(y)\,dy
   =\frac{L_t(\eta)}{1+2\cos(h\eta+\pi/3)}.
\]
Applying the same sech integral and derivative at the shifts
$2a,a,0,-a$ shows that its inverse is
\begin{equation}\label{rhosource}
 \rho_t(y)=\frac{A_te^{2ay}+(-y-t+3/2)e^{ay}/\sqrt2+C_t
                +[(1-t)y/2-t+3/2]e^{-ay}/\sqrt2}{2\cosh(3ay)}.
\end{equation}
Fourier uniqueness identifies this expression with the convolution.
The identity
\[w_y=p\cosh(ay)/(2\cosh(3ay))\]
shows that its quotient by $w_y$ is
exactly \eqref{qsource}. At positive infinity
$\rho_t(y)=A_te^{-ay}+O((1+y)e^{-2ay})$; at negative infinity
$\rho_t(y)=O((1+|y|)e^{2ay})$. As
$w_y(y)\sim(y^2/2)e^{-2a|y|}$, one has
$|\rho_t|^2/w_y=O(y^{-2})$ on the positive tail and exponential
decay on the negative one. Thus $q_t\in\mathcal H$.
The basis change and finite inverse formula now prove
\eqref{boundaryexact}. The norm \eqref{Hnorm} also follows by multiplying
the exact recurrence
$H_j/H_{j-1}=3j(j+2)(3j+2)(3j+4)/[4(2j+1)(2j+3)]$, $H_0=2$.

Finally set $A_*=3\sqrt3/4=e^{2\alpha}$. Stirling in \eqref{Hnorm}
gives
$\sqrt{h_j/H_j}\sim [8/(9\,3^{1/4})]A_*^{-j}$.
Equation \eqref{calibrationequiv} gives
$cZ_n(c)/Z_{n+1}(c)\sim[c/(\sqrt2 A_*)]A_*^{-n}$.
Substitution in \eqref{boundaryexact}, with $d_{n-1}\sim n$,
proves \eqref{calcoefficient}, including its alternating sign.
\end{proof}

\subsection{Stable inverses and the analytic coefficient reduction}
For completeness, the finite nonnormal inverses are uniformly bounded
on compact subsets of $\Omega$. In the unweighted representation,
$\widetilde{\mathcal B}=C_0+\mathcal E$, and the earlier proof gives
\[
 \delta(I+\delta P_nC_0P_n)^{-1}P_n\mathcal EP_n
 \longrightarrow\delta(I+\delta C_0)^{-1}\mathcal E
\]
in Hilbert--Schmidt norm, locally uniformly. Nonvanishing of
$\mathcal A$ makes $I+\delta(I+\delta C_0)^{-1}\mathcal E$
invertible. The exact factorization of $I+\delta P_n\widetilde{\mathcal B}P_n$
then bounds its inverse for all sufficiently large $n$; each of the
remaining finitely many sizes is controlled by the negative-root
property. Taking adjoints gives the corresponding bound for
$\mathcal B^*$. The resulting strong convergence on fixed vectors
still does not control a coefficient after division by
$a_{n-1}^{\rm src}=O(1/n)$.

Every $F_n$ is holomorphic and nowhere zero on $\HH_+$, with
$F_n(c)=1$. The earlier interlacing argument gives, for
$R_n(s)=\Log[Z_n(s)/Z_n(c)]$,
\[
 |R_{n+1}(s)-R_n(s)|\le
             \frac{\pi|s-c|}{\sqrt{c\Re s}}.
\]
Thus $F_n$ and $1/F_n$ are locally uniformly bounded. A locally
convergent subsequence and all its derivatives converge by Cauchy's
formula. If every fixed Taylor coefficient at $t=3$ of
$\gamma_{n-1}(t)/a_{n-1}^{\rm src}$ converges, every subsequential
analytic limit has the same Taylor series at $s=c$. Analytic
uniqueness and normal-family compactness then prove full-sequence
locally uniform convergence on $\HH_+$.

For $k\ge1$ the relevant coefficient is exactly
\begin{equation}\label{Taylorcoef}
 \frac{(-1)^k}{a_{n-1}^{\rm src}}\left\{
 \langle\psi_{n-1},(P_n\mathcal B^*P_n)^kP_nq_3\rangle
 -\langle\psi_{n-1},(P_n\mathcal B^*P_n)^{k-1}P_nq^{(1)}\rangle
 \right\}.
\end{equation}
This is the finite Neumann derivative at zero; it involves no
interchange of an infinite series and a size limit. We prove
existence of each fixed-power limit next. The established companion
limit will then evaluate the analytic boundary limit without any
need to compute those powers separately.
